% Part: incompleteness
% Chapter: representability-in-q
% Section: beta-function

\documentclass[../../../include/open-logic-section]{subfiles}

\begin{document}

\olfileid{inc}{req}{bet}
\olsection{لمة دالة بيتا}


لكي نبيّن إمكان تنفيذ العودية البدائية عند توافر الجمع والضرب
و$\Char{=}$، نحتاج إلى تطوير دوال تتعامل مع المتتاليات. (ولو كان
الرفع إلى قوة متاحًا أيضًا لكانت مهمتنا أسهل.) فعندما كانت العودية
البدائية متاحة، أمكننا تعريف أشياء مثل «العدد الأولي ذي الرتبة $n$»
واختيار ترميز مباشر نسبيًا. أما هنا فلا تتوافر لنا العودية
البدائية---بل نريد أن نثبت إمكان تنفيذها باستعمال التصغير---ولذلك
علينا أن نكون أكثر حيلة.

\begin{lem}
\ollabel{lem:beta}
توجد دالة $\beta(d,i)$ بحيث توجد، لكل متتالية $a_0$، \dots،~$a_n$،
قيمة~$d$ تحقق $\beta(d,i)=a_i$ لكل $i\le n$. وفوق ذلك يمكن تعريف
$\beta$ من الدوال الأساسية باستعمال التركيب والتصغير المنتظم وحدهما.
\end{lem}

اعتبر $d$ ترميزًا للمتتالية $\tuple{a_0,\dots,a_n}$، واعتبر
$\beta(d,i)$ معيدةً للعنصر ذي الرتبة $i$. (لاحظ أن هذا «الترميز»
\emph{لا} يستعمل ترميز قوى الأعداد الأولية الذي ألفناه!) واللمة
مقتصدة جدًا؛ فهي لا تقول إننا نستطيع وصل المتتاليات أو إلحاق العناصر،
ولا حتى إننا نستطيع \emph{حساب}~$d$ من $a_0$، \dots،~$a_n$ باستعمال
دوال قابلة للتعريف بالتركيب والتصغير المنتظم. وكل ما تقوله هو وجود دالة
«فك ترميز» تجعل كل متتالية «مرمّزة».

استعمال الرمز $\beta$ من وضع غودل. ونكرر أن الجزء الصعب من برهان اللمة
هو تعريف~$\beta$ ملائمة باستعمال الموارد المقيدة ظاهريًا، أي التركيب
والتصغير وحدهما---مع السماح لنا باستعمال الجمع والضرب و$\Char{=}$. ثمة
طرائق شتى لبرهان هذه اللمة، لكن طريقة غودل الأصلية لا تزال من أنظفها؛
وهي تستعمل حقيقة في نظرية الأعداد تسمى مبرهنة سون تزو (وتسمى تقليديًا
«مبرهنة الباقي الصينية»).

\begin{defn}
يكون العددان الطبيعيان $a$ و$b$ \emph{أوليين فيما بينهما} إذا وفقط
إذا كان قاسمهما المشترك الأكبر~$1$؛ أي لا يشتركان في قواسم أخرى.
\end{defn}

\begin{defn}
يكون العددان الطبيعيان $a$ و$b$ \emph{متطابقين بترديد~$c$}،
$a\equiv b\mod c$، إذا وفقط إذا كان $c\mid(a-b)$؛ أي إذا كان لهما
الباقي نفسه عند القسمة على~$c$.
\end{defn}

هذه هي مبرهنة سون تزو:
\begin{thm}
افترض أن $x_0$، \dots،~$x_n$ أولية فيما بينها مثنى مثنى، ولتكن
$y_0$، \dots،~$y_n$ أعدادًا كيفما كانت. عندئذ يوجد عدد $z$ بحيث
\begin{align*}
z & \equiv y_0 \mod x_0 \\
z & \equiv y_1 \mod x_1 \\
& \vdots  \\
z & \equiv y_n \mod x_n.
\end{align*}
\end{thm}

سنستعمل مبرهنة سون تزو كما يأتي: إذا كانت $x_0$، \dots،~$x_n$ أكبر
من $y_0$، \dots،~$y_n$ على الترتيب، أمكننا اتخاذ $z$ شفرةً للمتتالية
$\tuple{y_0,\dots,y_n}$. ولاستعادة~$y_i$ لا يلزم إلا قسمة $z$ على
$x_i$ وأخذ الباقي. ولاستخدام هذا الترميز سنحتاج إلى إيجاد قيم ملائمة
لـ$x_0$، \dots،~$x_n$.

ستفيدنا في ذلك ملاحظتان. إذا أُعطيت $y_0$، \dots،~$y_n$، فضع
\begin{align*}
j &= \max(n, y_0 + 1, \dots, y_n + 1), \\
m &= \lcm(1,\dots,j),
\end{align*}
وضع
\begin{align*}
x_0 & = 1 + m \\
x_1 & = 1 + 2 \cdot m \\
x_2 & = 1 + 3 \cdot m \\
& \vdots  \\
x_n & = 1 + (n+1) \cdot m
\end{align*}
عندئذ يصدق أمران:
\begin{enumerate}
\item\ollabel{rel-prime} الأعداد $x_0,\dots,x_n$ أولية فيما بينها.
\item\ollabel{less} لكل $i$، يكون $y_i<x_i$.
\end{enumerate}
ولرؤية صدق \olref{rel-prime}، لاحظ أنه إذا كان $p$ عددًا أوليًا وكان
$p\mid x_i$ و$p\mid x_k$، فإن $p\mid1+(i+1)m$ و
$p\mid1+(k+1)m$. لكن $p$ يقسم عندئذ فرقهما
\[
(1 + (i+1)m) - (1+ (k+1)m) = (i-k) m.
\]
ولما كان $p$ يقسم $1+(i+1)m$، فلا يمكن أن يقسم $m$ أيضًا (وإلا تركت
القسمة الأولى باقيًا قدره~$1$). ولذلك يقسم $p$ العدد $i-k$، لأنه يقسم
$(i-k)m$. لكن $\left|i-k\right|$ لا يتجاوز~$n$، وقد اخترنا
$j\geq n$، فيلزم من ذلك أن $p\mid m$، وهو تناقض مرة أخرى. ومن ثم لا
يوجد عدد أولي يقسم كلًا من $x_i$ و$x_k$. أما البند~\olref{less} فسهل:
لدينا $y_i<j\leq m<x_i$.

لنبرهن الآن لمة دالة $\beta$. تذكّر أنه يجوز لنا استعمال $0$ والخلف
والجمع والضرب و$\Char{=}$ والإسقاطات، وأية دالة معرّفة منها باستعمال
التركيب والتصغير المطبق على دوال منتظمة. ويجوز لنا أيضًا استعمال علاقة
إذا أمكن تعريف دالتها المميزة بهذه الطريقة. ويمكننا، كما سبق، أن نثبت
انغلاق هذه العلاقات تحت التركيبات البوليانية والتكميم المحدود؛ فمثلًا:
\begin{align*}
\fn{not}(x) & \defis \Char{=}(x,0)\\
\bmin{x \leq z}{R(x,y)} & \defis \umin{x}{(R(x,y) \lor x = z)}\\
\bexists{x \leq z}{R(x,y)} & \defiff R(\bmin{x \leq z}{R(x,y)}, y)
\end{align*}
ويمكننا بعد ذلك أن نثبت أن جميع الآتي قابل للتعريف أيضًا من دون عودية
بدائية:
\begin{enumerate}
\item دالة المزاوجة، $J(x,y)=\frac{1}{2}[(x+y)(x+y+1)]+x$؛
% maybe explain more what is going on here, a bit confusing.
\item دالتا الإسقاط
\begin{align*}
K(z) & = \bmin{x \leq z}{\bexists{y \leq z}{z = J(x,y)}},\\
L(z) & = \bmin{y \leq z}{\bexists{x \leq z}{z = J(x,y)}};
\end{align*}
\item علاقة الأصغر $x<y$؛
\item علاقة القسمة $x\mid y$؛
% \item $x \tsub y$
% \item $\fn{Prime}(x)$
% \item Assuming $p$ is prime, the relation ``$x$ is a power of $p$'':
% \[
% \bforall{y \leq x}{(y \mid x \lif y = 1 \lor y = x)}.
% \]
\item الدالة $\fn{rem}(x,y)$ التي تعيد باقي قسمة $y$ على~$x$.
\end{enumerate}
والآن عرّف
\begin{align*}
\beta^*(d_0,d_1,i) & = \fn{rem}(1+(i+1) d_1,d_0) \text{ و}\\
\beta(d,i) & = \beta^*(K(d),L(d),i).
\end{align*}
هذه هي الدالة المطلوبة. إذا أُعطيت $a_0,\dots,a_n$ كما أعلاه، فضع
\[
j = \max(n,a_0+1,\dots,a_n+1),
\]
ولتكن $d_1=\lcm(1,\dots,j)$. وبحسب \olref{rel-prime} أعلاه، نعلم أن
$1+d_1$ و$1+2d_1$، \dots، $1+(n+1)d_1$ أولية فيما بينها، وبحسب
\olref{less} فإنها أكبر من $a_0,\dots,a_n$ على الترتيب. وبمبرهنة سون
تزو توجد قيمة~$d_0$ بحيث، لكل~$i$،
\[
d_0 \equiv a_i \mod (1+(i+1)d_1)
\]
ومن ثم (لأن $d_1$ أكبر من~$a_i$)،
\[
a_i = \fn{rem}(1+(i+1)d_1,d_0).
\]
ضع $d=J(d_0,d_1)$. عندئذ، لكل $i\le n$، يكون
\begin{align*}
\beta(d,i) & =  \beta^*(d_0,d_1,i) \\
& =  \fn{rem}(1+(i+1) d_1,d_0) \\
& =  a_i
\end{align*}
وهذا هو المطلوب. وبه يتم برهان لمة دالة~$\beta$.

\begin{prob}
  أثبت أن العلاقتين $x<y$ و$x\mid y$ والدالة~$\fn{rem}(x,y)$ قابلة
  للتعريف من دون عودية بدائية. ويجوز لك استعمال $0$ والخلف والجمع
  والضرب و$\Char{=}$ والإسقاطات والتصغير والتكميم المحدودين.
\end{prob}

\end{document}
